\documentclass{jarticle}
\usepackage[dvipdfmx]{graphicx}
\topmargin       -10.0mm
\oddsidemargin   -0.4mm
\evensidemargin  -0.4mm
\headheight       0.0mm
\headsep          0.0mm
\footskip        10.0mm
%
\textheight     260.0mm
\textwidth      170.0mm
%
\setlength{\parindent}{1zw}
%
\usepackage{graphicx}
\usepackage{multirow}
\usepackage{wrapfig}
\makeatletter
\def \fgcaption{\def \@captype{figure} \caption}
\def \tbcaption{\def \@captype{table} \caption}
\makeatother
\def \figurename{Fig.}
\def \tablename{Table}
%
\def\topfraction{.99}
\def\bottomfraction{.99}
\def\textfraction{.01}
\def\floatpagefraction{.99}
%
\pagestyle{empty}
%
\begin{document}
%
\begin{center}
 \begin{tabular}{|c|c|c|} \hline
  \parbox{50mm}{\vspace{1mm} \centering{\Large\bf 2026/06/04}} & 
  \parbox{55mm}{\vspace{1mm} \centering{\Large\bf B4　堀田裕樹}} & 
  \parbox{50mm}{\vspace{1mm} \centering{\Large\bf No.~8}}
  \\ [2mm] \hline
%
  \multicolumn{3}{|c|}{%
   \parbox{155mm}{\centering\Large\bf
    \vspace{1mm}
    シミュレーションデータの解析条件の確立と特定パラメータの比較評価
   }
  }
  \\ [2mm] \hline\hline
%
  \multicolumn{3}{|c|}{%
   \parbox{155mm}{\large\vspace{1mm}
%
    {\Large\bf 今月の目標}
%
    \begin{itemize}
     \item シミュレーションの実行
     \item シミュレーションデータの処理
    \end{itemize}
%
    {\Large\bf 進捗の要点}
%
    \begin{itemize}
     \item 先行研究に基づき，波形切り出し範囲をピーク前後512点（計1024点）に設定
     \item FFTスペクトルの評価周波数範囲を 2--10 MHz に確定
    \end{itemize}
%
    {\Large\bf 今後の実施事項}
%
    \begin{itemize}
     \item 実験及びそのデータ処理の実施
    \end{itemize}
%
   }
  }
  \\ [2mm] \hline
 \end{tabular}
%
\end{center}
%
\par
本週は，具体的な解析条件の確立とパラメータ比較を行った．
解析の妥当性確保のため，波形切り出し範囲を先行研究と統一し，最大ピーク前後512点（計1024点）とした．FFTスペクトルは，探触子特性を考慮し 2--10 MHz を解析帯域とした．

これらを用い格子幅(width)の比較を行った．Fig.\ref{fig:results} に Pitch: 2000 および Pitch: 1250 の波形・FFT結果を示す．width が小さい場合、反射波の振幅が大きくテイル波の振幅は小さくなることが確認された.
また，FFTスペクトルにおいては，幅が大きい場合、ディップ周波数が高くなることが確認された．
\begin{table}[htb]
  \centering
  \tbcaption{Simulation parameters (Unit: $10^{-5}$m)}
  \label{tab:params}
  \begin{tabular}{ccccc} \hline
    Parameter & \multicolumn{4}{c}{Values} \\ \hline
    Pitch & 2000 & 1750 & 1500 & 1250 \\
    Depth & 200 & 150 & 125 & 100 \\
    Width & 800 & 600 & 400 & 150 \\ \hline
  \end{tabular}
\end{table}

\begin{center}
  \begin{minipage}[b]{0.49\linewidth}
    \centering
    \includegraphics[width=6.8cm]{../../pic/report_20260604/p2000_waveform.png}
    \vspace{-2mm}
    \fgcaption{Waveform (Pitch: 2000 Depth: 200)}
  \end{minipage}
  \begin{minipage}[b]{0.49\linewidth}
    \centering
    \includegraphics[width=6.8cm]{../../pic/report_20260604/p2000_fft.png}
    \vspace{-2mm}
    \fgcaption{FFT (Pitch: 2000 Depth: 200)}
  \end{minipage}
  \\ \vspace{2mm}
  \begin{minipage}[b]{0.49\linewidth}
    \centering
    \includegraphics[width=6.8cm]{../../pic/report_20260604/p1250_waveform.png}
    \vspace{-2mm}
    \fgcaption{Waveform (Pitch: 1250 Depth: 200)}
  \end{minipage}
  \begin{minipage}[b]{0.49\linewidth}
    \centering
    \includegraphics[width=6.8cm]{../../pic/report_20260604/p1250_fft.png}
    \vspace{-2mm}
    \fgcaption{FFT (Pitch: 1250 Depth: 200)}
  \end{minipage}
  \label{fig:results}
\end{center}

\vspace{-6mm}
\begin{thebibliography}{9}

\end{thebibliography}
%
\end{document}
